\textbf{Llano no es una herramienta para ahorrar tokens.} Frente a la instrucción de concisión, el largo no cambia: \deltaWordsTerse{} en palabras y \deltaTokensTerse{} en tokens visibles, con intervalos que incluyen el cero. Toda la reducción frente a la respuesta sin instrucción (\deltaWordsBaseline{} en palabras) la consigue ya la frase ``Responde de forma concisa''. Caveman produce respuestas más cortas que llano: llano usa \deltaWordsCaveman{} más palabras. Este resultado coincide con la lección del harness de caveman: casi todo el ahorro de un estilo viene del pedido genérico de brevedad.

\textbf{El efecto de llano está en la forma.} Las violaciones cada 100 palabras bajan de \tersePerHundred{} a \llanoPerHundred. Llano tiene menos violaciones que la instrucción de concisión en \violBetterTerse{} prompts, las mismas en \violEqualTerse{} y más en \violWorseTerse{} ($p$ \violPTerse{} en una prueba de signos). Las frases largas, el punto y coma y las muletillas casi desaparecen. Contra caveman la diferencia también es clara: \violBetterCaveman{} prompts mejor, \violWorseCaveman{} peor ($p = \violPCaveman$). La legibilidad INFLESZ sube de \terseInflesz{} a \llanoInflesz{}, de ``bastante fácil'' hacia el borde de ``muy fácil''.

Esta métrica tiene un sesgo a favor de llano: el linter mide las reglas que llano conoce. Por eso importan las dos métricas siguientes.

\textbf{Llano no pierde información.} La fidelidad es \llanoFidelity{} con llano, \terseFidelity{} con la instrucción de concisión y \baselineFidelity{} sin instrucción. Los hechos que faltan son casi siempre detalles secundarios, como la versión desplegada o el paso exacto del Dockerfile. Las afirmaciones incorrectas son pocas en todos los grupos. Caveman tiene el doble que llano.

\textbf{La preferencia a ciegas favorece a llano, pero sin significancia.} El juez prefirió llano en \winsTerse{} de \evalPrompts{} comparaciones contra la instrucción de concisión y en \winsCaveman{} contra caveman ($p = \pSignTerse$ en ambos casos). Contra la respuesta sin instrucción ganó \winsBaseline{} veces ($p = \pSignBaseline$). Con \evalPrompts{} prompts y una sola corrida, el piloto no alcanza para confirmar esta preferencia.

\textbf{El costo es la entrada.} Llano agrega su skill y el paquete de idioma a cada llamada: unos \llanoInputTokens{} tokens contra \terseInputTokens{}. Con la caché de prompts, ese costo baja mucho en una sesión real, pero el texto sigue ocupando contexto. Reducir el paquete, por ejemplo cargando el diccionario solo en modo \texttt{all}, es la mejora más directa.

\textbf{Próximos pasos.} Ejecutar tres corridas con Opus, Sonnet y Haiku para medir la variación y evitar que el juez evalúe solo a su propio modelo. Ampliar el conjunto de prompts con casos escritos por otras personas. Medir una versión de llano con un paquete más corto.
